%===============================================================================
\section{Kết luận điều hành (Executive Summary)}
%===============================================================================

Tổng thể mã nguồn xác nhận repository là một \textbf{nền tảng MLOps PaaS mạnh mẽ ở khâu vận hành và triển khai mô hình (Model Operations)}, nhưng chưa hoàn thiện năng lực \textbf{tự thích ứng và bảo trì mô hình khép kín (Closed-Loop Continuous Training)}. Điểm mạnh nổi bật nằm ở kiến trúc phân tầng rõ ràng (Control Plane vs. Data Plane vs. Execution Plane), tính bất biến của phiên bản mô hình, tính đa dạng của backend huấn luyện và kênh thu thập sự kiện suy luận. Ngược lại, điểm nghẽn lớn nhất nằm ở sự thiếu vắng vòng lặp phản hồi nhãn thực tế, thiếu cơ chế giám sát chất lượng/hiệu năng suy giảm và chưa có trí tuệ ra quyết định (decision intelligence) giữa khâu cảnh báo và kích hoạt tái huấn luyện.

\begin{table}[ht]
\centering
\caption{Đánh giá tổng hợp năng lực hệ thống theo khung kiến trúc tham chiếu.}
\begin{tabularx}{\linewidth}{@{}L{4.4cm} C{1.6cm} Y@{}}
\toprule
\textbf{Trục kiến trúc} & \textbf{Mức độ} & \textbf{Nhận định tổng quan} \\
\midrule
Nền tảng \& Phân tầng & 7/10 & Phân tách rõ ràng giữa Control, Data và Execution Plane; phân quyền tenant chuẩn mực. \\
Lưu trữ \& Quản lý phiên bản & 7/10 & Quản lý artifact, digest image và snapshot huấn luyện bất biến; thiếu Feedback DB. \\
Huấn luyện \& CI/CD & 7/10 & Hỗ trợ đa backend (Docker/Argo/Kubeflow), đóng gói tự động; thiếu Evaluation Gate tự động. \\
Suy luận \& Quan sát & 6/10 & Gateway serving tốt, thu thập event stream ổn định; chưa đo lường performance thực tế. \\
Vòng đời bảo trì khép kín & 4/10 & Tự động kích hoạt DriftRun khi đủ mẫu; chưa có cơ chế chẩn đoán nguyên nhân và CT policy. \\
Quản trị \& Stakeholder & 4/10 & Có xác thực JWT/API-Key và audit log; chưa có luồng phê duyệt (approval) cho bảo trì mô hình. \\
\bottomrule
\end{tabularx}
\end{table}

\textbf{Xác định biến thể kiến trúc hiện tại:} Đối chiếu với phân loại của Amou Najafabadi \textit{et al.}~\cite{najafabadi2026}, hệ thống hiện đáp ứng trọn vẹn đặc trưng của \textbf{Biến thể V1 (Inference Service + Manual Retraining)}: dịch vụ \codepath{model-packager} đóng gói mô hình thành container độc lập, task \codepath{deployment} đẩy lên hạ tầng production, và \codepath{model-server} định tuyến phục vụ dự đoán. Việc tái huấn luyện hoàn toàn phụ thuộc vào thao tác thủ công của kỹ sư ML. Mục tiêu của đề tài là nghiên cứu và nâng cấp hệ thống từ biến thể V1 lên \textbf{Biến thể V3 (Continuous Training Pipeline)} có tích hợp cơ chế nhận thức suy giảm thích ứng.
